\section{期待と異なるデータを検証する}
\label{theory:research:validation}
測定結果が仮説や依頼者の期待に反したときには、期待からの距離をデータの品質の尺度にしない。装置、校正、測定範囲、試料、時刻、記録の処理を点検し、除外には科学的な理由を記録する。予測値と実測値は区別する。測定できなかった部分を予測で補う場合にも、観測として扱わず方法と限界を示す\parencite{scj_code2013,mext_misconduct2014}。

報告期限と装置故障が重なる場合は、修理や代替測定の見通しを調べ、指導者・依頼者と期限、暫定報告、再測定を協議する。すべての未確認値を無説明で提出することも、都合の悪い値を無根拠に捨てることも妥当でない。生データと判断の根拠を残すことで、結論を後から検証できる。

\paragraph{説明と記録を分けて管理する}
依頼者が詳しくないことは、重要な仕様変更や危険を説明しなくてよい根拠にはならない。リスクを相手の判断に使える形で示す。技術情報を他者へ伝える場合は公開性、契約、守秘範囲を確かめる。すべての技術的会話が秘密漏洩ではない一方、業務上の非公開情報を友人へ話してよいともしない。

測定記録の保存・廃棄には責任、期間、アクセスの管理が必要である。記憶に基づく説明がすべて不正というわけではないが、記録があるのに確認せず数値報告を作れば正確性と検証可能性を失う。パスワードや施錠は対策の一部であり、更新・権限・保存期間も点検する。

\section{研究不正への個人責任と機関責任}
\label{theory:research:guideline2014}
平成26年の研究不正ガイドラインは、個人の不正の認定と措置に加え、研究機関の教育・予防・調査の責任を扱う。機関への管理条件、間接経費の削減、配分停止等の段階があり、「行政措置の根拠にはならない」とする説明は適切でない\parencite{mext_misconduct2014}。

教育対象を研究者本人だけに狭めない。学生や研究支援者など、実際に研究へ関わる人々に対する説明と組織体制を考える。研究活動の不正、研究費の不正使用、日常の不適切行為には異なる規程や手続がある。出題年のガイドラインの範囲・条件を確認し、現行の改定や別の資金制度の例外を一括して遡及しない。
